% Joint-paper experimental protocol for reward construction.
\subsection{Reward Construction on LIBERO}\label{sec:libero_setup}
We test whether task-progress feedback benefits a trained policy when complete-task outcomes provide limited contrast, using the phase-aware reward in Sec.~\ref{sec:dense_method}. These experiments evaluate reward construction with an established action interface.

We train OpenVLA-OFT in SimpleVLA-RL on the ten LIBERO-Spatial tasks, using two NVIDIA L40 GPUs and an aggregate rollout batch of 32. LoRA rank and alpha are both 16, applied to the language-model projector; gradient checkpointing and FSDP offloading reduce memory demand. These updates optimize discrete action-token likelihoods, despite the continuous-action supervised formulation associated with OFT. \pending{Confirm implementation commit, tokenizer/action decoding, adapted modules, group size, and optimization schedule.}

All aggregate reward configurations start from 61.0\% validation success. The primary dense condition uses $(w_p,w_c,C,\lambda)=(0.2,0.3,0.05,5)$; its terminal-only control disables shaping and retains $\lambda=5$. The filtered dense condition additionally uses $[L,U]=[0.1,0.9]$. With $\lambda=1$, dense shaping uses $w_p=w_c=0.05$, while phase-only shaping uses $w_p=0,w_c=0.2$; both use $C=0.05$.

We report the final recorded evaluations, balanced across the ten tasks. To examine low-success learning, we restrict training and validation to \texttt{task\_1}, starting at 10\%, with $[L,U]=[0.1,1.0]$ in both conditions. \pending{Complete evaluation counts, seeds, canonical task name, elapsed budgets, and accepted-update counts.} Percentages are descriptive until these records support uncertainty estimates and budget matching.
